\chapter{2008年江西卷（理）}

\section{单选题}

\begin{problem}
在复平面内，复数 \(z = \sin 2 + \i\cos 2\) 对应的点位于（\quad）

\choices
    {第一象限}
    {第二象限}
    {第三象限}
    {第四象限}
\end{problem}

\begin{answer}
D
\end{answer}

\begin{solution}
因为 \(\frac{\pi}{2}<2<\pi\)，所以 \(\sin 2>0\)，\(\cos 2<0\). 复数 \(z=\sin 2+\i\cos 2\) 对应的点的横坐标为正、纵坐标为负，故该点位于第四象限.
\end{solution}

\begin{problem}
定义集合运算： \( A * B = \{z \mid z = xy, x \in A, y \in B\} \). 设 \( A = \{1, 2\} \)，\( B = \{0, 2\} \)，则集合 \( A * B \) 的所有元素之和为（\quad）

\choices
    {\(0\)}
    {\(2\)}
    {\(3\)}
    {\(6\)}
\end{problem}

\begin{answer}
D
\end{answer}

\begin{solution}
根据定义，取 \(x\in A\)、\(y\in B\) 可得乘积 \(1\cdot0=0\)、\(1\cdot2=2\)、\(2\cdot0=0\)、\(2\cdot2=4\). 因此
\[
A*B=\{0,2,4\}
\]
其所有元素之和为 \(0+2+4=6\)，故选 D.
\end{solution}

\begin{problem}
若函数 \( y = f(x) \) 的值域是 \( \left[\frac{1}{2}, 3\right] \)，则函数 \( F(x) = f(x) + \frac{1}{f(x)} \) 的值域是（\quad）

\choices
    {\(\left[\frac{1}{2}, 3\right]\)}
    {\(\left[2, \frac{10}{3}\right]\)}
    {\( \left[\frac{5}{2}, \frac{10}{3}\right] \)}
    {\(\left[3, \frac{10}{3}\right]\)}
\end{problem}

\begin{answer}
B
\end{answer}

\begin{solution}
令 \(t=f(x)\)，则 \(t\in\left[\frac12,3\right]\)，且所求函数值为
\[
F=t+\frac1t
\]
在 \(t>0\) 时，函数 \(t+\frac1t\) 的导数为 \(1-\frac1{t^2}\)，所以它在 \(\left[\frac12,1\right]\) 上递减，在 \([1,3]\) 上递增. 因而最小值为 \(2\)，最大值为
\[
\max\left\{\frac12+2,\,3+\frac13\right\}=\frac{10}{3}
\]
所以值域为 \(\left[2,\frac{10}{3}\right]\)，故选 B.
\end{solution}

\begin{problem}
\(\displaystyle \lim\limits_{x\to 1}\frac{\sqrt{x + 3} - 2}{\sqrt{x} - 1} =\) （\quad）

\choices
    {\( \frac{1}{2} \)}
    {\(0\)}
    {\(-\frac{1}{2}\)}
    {不存在}
\end{problem}

\begin{answer}
A
\end{answer}

\begin{solution}
当 \(x\ne1\) 时，分子、分母分别有理化，得
\[
\frac{\sqrt{x+3}-2}{\sqrt{x}-1}
=\frac{x-1}{\sqrt{x+3}+2}\cdot\frac{\sqrt{x}+1}{x-1}
=\frac{\sqrt{x}+1}{\sqrt{x+3}+2}
\]
令 \(x\to1\)，得到
\[
\lim_{x\to1}\frac{\sqrt{x+3}-2}{\sqrt{x}-1}=\frac{2}{4}=\frac12
\]
故选 A.
\end{solution}

\begin{problem}
在数列 \(\{a_{n}\}\) 中，\(a_{1} = 2\)，\(a_{n+1} = a_{n} + \ln \left(1 + \frac{1}{n}\right)\)，则 \(a_{n} =\)（\quad）

\choices
    {\(2 + \ln n\)}
    {\(2 + (n - 1)\ln n\)}
    {\(2 + n\ln n\)}
    {\(1 + n + \ln n\)}
\end{problem}

\begin{answer}
A
\end{answer}

\begin{solution}
由递推式逐项相加，得到
\[
a_n=a_1+\sum_{k=1}^{n-1}\ln\left(1+\frac1k\right)
=2+\sum_{k=1}^{n-1}\ln\frac{k+1}{k}
\]
利用对数的运算性质，
\[
\sum_{k=1}^{n-1}\ln\frac{k+1}{k}
=\ln\prod_{k=1}^{n-1}\frac{k+1}{k}=\ln n
\]
因此 \(a_n=2+\ln n\)，故选 A.
\end{solution}

\begin{problem}
函数 \( y = \tan x + \sin x - |\tan x - \sin x| \) 在区间 \(\left(\frac{\pi}{2}, \frac{3\pi}{2}\right)\) 内的图象大致是 （\quad）

\choices
    {\choicebitmap[width=0.15\paperwidth]{img/2008/jiangxi_science/q6_optA_nolabel.png}}
    {\choicebitmap[width=0.15\paperwidth]{img/2008/jiangxi_science/q6_optB_nolabel.png}}
    {\choicebitmap[width=0.15\paperwidth]{img/2008/jiangxi_science/q6_optC_nolabel.png}}
    {\choicebitmap[width=0.15\paperwidth]{img/2008/jiangxi_science/q6_optD_nolabel.png}}
\end{problem}

\begin{answer}
D
\end{answer}

\begin{solution}
在 \(\left(\frac\pi2,\pi\right)\) 内，\(\sin x>0\)、\(\cos x<0\)，所以 \(\tan x<0<\sin x\)，从而
\[
y=\tan x+\sin x-(\sin x-\tan x)=2\tan x
\]
该段图象从 \(x\to\frac\pi2+\) 时的负无穷上升到 \(x\to\pi-\) 时的 \(0\).

在 \(\left(\pi,\frac{3\pi}{2}\right)\) 内，\(\sin x<0\)、\(\cos x<0\)，所以 \(\tan x>0>\sin x\)，从而
\[
y=\tan x+\sin x-(\tan x-\sin x)=2\sin x
\]
在 \(x=\pi\) 时，\(y=0\)，该点属于定义域；第二段图象从 \(x\to\pi+\) 时的 \(0\) 下降到 \(x\to\frac{3\pi}{2}-\) 时的 \(-2\)，而 \(x=\frac\pi2\) 和 \(x=\frac{3\pi}{2}\) 不属于定义域. 这与选项 D 的图象相符，故选 D.
\end{solution}

\begin{problem}
已知 \(F_{1}\)，\(F_{2}\) 是椭圆的两个焦点，满足 \(\overrightarrow{MF_{1}} \cdot \overrightarrow{MF_{2}} = 0\) 的点 \(M\) 总在椭圆内部，则椭圆离心率的取值范围是（\quad）

\choices
    {\((0,1)\)}
    {\(\left(0, \frac{1}{2}\right]\)}
    {\(\left(0, \frac{\sqrt{2}}{2}\right)\)}
    {\(\left[\frac{\sqrt{2}}{2}, 1\right)\)}
\end{problem}

\begin{answer}
C
\end{answer}

\begin{solution}
设椭圆的半长轴、半短轴和半焦距分别为 \(a\)、\(b\)、\(c\)，则 \(a>b>0\)，且
\[
a^2=b^2+c^2
\]
取椭圆中心为原点、焦点为 \(F_1(-c,0)\)、\(F_2(c,0)\). 对点 \(M(x,y)\)，条件
\[
\overrightarrow{MF_1}\cdot\overrightarrow{MF_2}=0
\]
化为
\[
(x+c)(x-c)+y^2=0
\]
即 \(x^2+y^2=c^2\). 因而这些点构成以椭圆中心为圆心、半径为 \(c\) 的圆.

该圆上的点都在椭圆内部时，圆在短轴端点处最容易到达椭圆边界，因此必须有 \(c<b\). 于是
\[
e=\frac ca<\frac ba=\sqrt{1-e^2}
\]
由于 \(e>0\)，平方后得 \(2e^2<1\)，即
\[
0<e<\frac{\sqrt2}{2}
\]
故选 C.
\end{solution}

\begin{problem}
\((1 + \sqrt[3]{x})^6\left(1 + \frac{1}{\sqrt[4]{x}}\right)^{10}\) 展开式中的常数项为（\quad）

\choices
    {\(1\)}
    {\(46\)}
    {\(4245\)}
    {\(4246\)}
\end{problem}

\begin{answer}
D
\end{answer}

\begin{solution}
第一项取 \(\mathrm C_6^r(\sqrt[3]{x})^r\)，第二项取 \(\mathrm C_{10}^k(x^{-1/4})^k\)，所得项的幂为
\[
\frac r3-\frac k4
\]
成为常数项的条件是 \(\frac r3-\frac k4=0\)，即 \(4r=3k\). 在 \(0\leqslant r\leqslant6\)、\(0\leqslant k\leqslant10\) 内，只有
\((r,k)=(0,0)\)，\((3,4)\)，\((6,8)\).
所以常数项为
\[
\mathrm C_6^0\mathrm C_{10}^0+\mathrm C_6^3\mathrm C_{10}^4+\mathrm C_6^6\mathrm C_{10}^8
=1+20\times210+1\times45=4246
\]
故选 D.
\end{solution}

\begin{problem}
若 \(0 < a_{1} < a_{2}\)，\(0 < b_{1} < b_{2}\)，且 \(a_{1} + a_{2} = b_{1} + b_{2} = 1\)，则下列代数式中值最大的是（\quad）

\choices
    {\(a_{1}b_{1} + a_{2}b_{2}\)}
    {\(a_{1}a_{2} + b_{1}b_{2}\)}
    {\(a_{1}b_{2} + a_{2}b_{1}\)}
    {\(\frac{1}{2}\)}
\end{problem}

\begin{answer}
A
\end{answer}

\begin{solution}
令 \(u=a_2-a_1\)、\(v=b_2-b_1\)，则 \(0<u<1\)、\(0<v<1\)，并且
\(a_1=\frac{1-u}{2}\)，\(a_2=\frac{1+u}{2}\)，\(b_1=\frac{1-v}{2}\)，\(b_2=\frac{1+v}{2}\).
记四个选项中的代数式依次为 \(X_1\)，\(X_2\)，\(X_3\)，\(X_4\). 直接代入得
\(X_1=\frac{1+uv}{2}\)，\(X_2=\frac{2-u^2-v^2}{4}\)，\(X_3=\frac{1-uv}{2}\)，\(X_4=\frac12\).
于是
\[
X_1-X_3=uv>0
\]
\[
X_1-X_4=\frac{uv}{2}>0
\]
以及
\[
X_1-X_2=\frac{(u+v)^2}{4}>0
\]
因此 \(X_1\) 最大，故选 A.
\end{solution}

\begin{problem}
连结球面上两点的线段称为球的弦. 半径为 \(4\) 的球的两条弦 \(AB\)，\(CD\) 的长度分别等于 \(2\sqrt{7}\)，\(4\sqrt{3}\)，\(M\)，\(N\) 分别为 \(AB\)，\(CD\) 的中点，每条弦的两端都在球面上运动，有下列四个命题：

\begin{circlelist}
\item 弦 \(AB\)，\(CD\) 可能相交于点 \(M\)；
\item 弦 \(AB\)，\(CD\) 可能相交于点 \(N\)；
\item \(MN\) 的最大值为 \(5\)；
\item \(MN\) 的最小值为 \(1\).
\end{circlelist}

\choices
    {\(1\) 个}
    {\(2\) 个}
    {\(3\) 个}
    {\(4\) 个}
\end{problem}

\begin{answer}
C
\end{answer}

\begin{solution}
设球心为 \(O\). 弦中点到球心的连线垂直于弦，因此
\(OM=\sqrt{4^2-\left(\frac{2\sqrt7}{2}\right)^2}=3\)，\(ON=\sqrt{4^2-\left(\frac{4\sqrt3}{2}\right)^2}=2\).

弦 \(AB\) 的所在直线上各点到 \(O\) 的距离不小于直线到 \(O\) 的距离 \(OM=3\)，而 \(ON=2\)，所以弦 \(AB\) 不可能经过 \(N\)，第二个命题错误.

弦 \(CD\) 所在直线到 \(O\) 的距离为 \(2\). 可以取一条经过 \(M\) 且到 \(O\) 的距离为 \(2\) 的直线；其上从垂足到 \(M\) 的距离为 \(\sqrt{OM^2-ON^2}=\sqrt5<2\sqrt3\)，故以该垂足为中点截取长度为 \(4\sqrt3\) 的弦时，\(M\) 在线段 \(CD\) 上. 因此第一个命题正确.

又由三角不等式，
\[
|OM-ON|\leqslant MN\leqslant OM+ON
\]
即 \(1\leqslant MN\leqslant5\). 当 \(M\)，\(O\)，\(N\) 共线且同向时取最小值 \(1\)，反向时取最大值 \(5\)，所以第三、第四个命题均正确. 真命题共有 3 个，故选 C.
\end{solution}

\begin{problem}
电子钟一天显示的时间是从 \(00{:}00\) 到 \(23{:}59\). 每一时刻都由四个数字组成. 则一天中任一时刻显示的四个数字之和为 \(23\) 的概率为（\quad）

\choices
    {\(\frac{1}{180}\)}
    {\(\frac{1}{288}\)}
    {\(\frac{1}{360}\)}
    {\(\frac{1}{480}\)}
\end{problem}

\begin{answer}
C
\end{answer}

\begin{solution}
一天共有 \(24\times60=1440\) 个时刻，且每个时刻等可能.

若小时的十位为 \(0\)，四个数字和为 \(23\) 只能是 \(09{:}59\). 若小时的十位为 \(1\)，需满足“小时个位、分钟十位、分钟个位”之和为 \(22\)，只有 \(18{:}59\)、\(19{:}49\)、\(19{:}58\). 小时的十位为 \(2\) 时，四个数字和的最大值为 \(2+3+5+9=19\)，不可能为 \(23\).

所以满足条件的时刻共有 4 个，所求概率为
\[
\frac4{1440}=\frac1{360}
\]
故选 C.
\end{solution}

\begin{problem}
已知函数 \( f(x) = 2mx^2 - 2(4 - m)x + 1 \)，\( g(x) = mx \)，若对于任一实数 \( x \)，\( f(x) \) 与 \( g(x) \) 的值至少有一个为正数，则实数 \( m \) 的取值范围是（\quad）

\choices
    {\((0, 2)\)}
    {\((0, 8)\)}
    {\((2, 8)\)}
    {\((-\infty, 0)\)}
\end{problem}

\begin{answer}
B
\end{answer}

\begin{solution}
若 \(m<0\)，当 \(x\to+\infty\) 时，\(f(x)\to-\infty\)，且 \(g(x)=mx<0\)，所以对充分大的正数 \(x\)，二者都不是正数，条件不成立. 若 \(m=0\)，取 \(x=1\) 得 \(f(1)=-7\)、\(g(1)=0\)，条件也不成立. 因此必须有 \(m>0\).

当 \(m>0\) 且 \(x\geqslant0\) 时，若 \(x>0\)，则 \(g(x)>0\)；若 \(x=0\)，则 \(f(0)=1>0\). 所以只需保证 \(x<0\) 时 \(f(x)>0\).

此时
\[
f(x)=2mx^2+(2m-8)x+1
\]
当 \(0<m\leqslant4\) 时，\(x<0\) 有
\[
f'(x)=4mx+2m-8<0
\]
所以 \(f(x)>f(0)=1>0\). 当 \(m>4\) 时，抛物线顶点横坐标为
\[
x_0=-\frac{2m-8}{4m}=\frac{4-m}{2m}<0
\]
且顶点处的最小值为
\[
f(x_0)=1-\frac{(2m-8)^2}{8m}=-\frac{(m-2)(m-8)}{2m}
\]
要使该最小值为正，必须且只需 \(4<m<8\). 结合前一种情形，得到
\[
0<m<8
\]
故选 B.
\end{solution}

\section{填空题}

\begin{problem}
直角坐标平面内三点 \(A(1,2)\)，\(B(3,-2)\)，\(C(9,7)\)，若 \(E\)，\(F\) 为线段 \(BC\) 的三等分点，则 \(\overrightarrow{AE} \cdot \overrightarrow{AF} = \)\fillinblank{}.
\end{problem}

\begin{answer}
22
\end{answer}

\begin{solution}
由 \(\overrightarrow{BC}=(6,9)\)，线段的两个三等分点为
\(E=B+\frac13\overrightarrow{BC}=(5,1)\)，\(F=B+\frac23\overrightarrow{BC}=(7,4)\).
所以
\(\overrightarrow{AE}=(4,-1)\)，\(\overrightarrow{AF}=(6,2)\)
从而
\[
\overrightarrow{AE}\cdot\overrightarrow{AF}=4\times6+(-1)\times2=22
\]
\end{solution}

\begin{problem}
不等式 \(2^{x - \frac{3}{x} + 1} \leqslant \frac{1}{2}\) 的解集为\fillinblank{}.
\end{problem}

\begin{answer}
\(\left(-\infty,-3\right]\cup\left(0,1\right]\)
\end{answer}

\begin{solution}
因为底数 \(2>1\)，原不等式等价于
\(x-\frac3x+1\leqslant-1\)，即 \(x-\frac3x+2\leqslant0\)，且 \(x\ne0\).
当 \(x>0\) 时，乘以正数 \(x\)，得
\[
x^2+2x-3=(x-1)(x+3)\leqslant0
\]
故 \(0<x\leqslant1\). 当 \(x<0\) 时，乘以负数 \(x\) 后不等号改变方向，得
\[
x^2+2x-3\geqslant0
\]
故 \(x\leqslant-3\). 因此解集为
\[
(-\infty,-3]\cup(0,1]
\]
\end{solution}

\begin{problem}
过抛物线 \(x^{2} = 2py\) （\(p > 0\)）的焦点 \(F\) 作倾斜角为 \(30^{\circ}\) 的直线，与抛物线分别交于 \(A\)，\(B\) 两点（点 \(A\) 在 \(y\) 轴左侧），则 \(\frac{|AF|}{|FB|} =\)\fillinblank{}.
\end{problem}

\begin{answer}
\(\frac{1}{3}\)
\end{answer}

\begin{solution}
抛物线 \(x^2=2py\) 的焦点为 \(F\left(0,\frac p2\right)\). 过焦点且倾斜角为 \(30^\circ\) 的直线方程为
\[
y=\frac{x}{\sqrt3}+\frac p2
\]
与抛物线联立，得
\(x^2=2p\left(\frac{x}{\sqrt3}+\frac p2\right)\)，\(x^2-\frac{2p}{\sqrt3}x-p^2=0\).
其两个根分别为
\(x_A=-\frac p{\sqrt3}\)，\(x_B=\sqrt3p\)
其中 \(x_A<0<x_B\). 因为 \(A\)，\(F\)，\(B\) 共线，且横坐标差与同一直线上距离成固定比例，所以
\[
\frac{|AF|}{|FB|}=\frac{|x_A-0|}{|x_B-0|}
=\frac{p/\sqrt3}{\sqrt3p}=\frac13
\]
\end{solution}

\begin{problem}
如图 1，一个正四棱柱形的密闭容器水平放置，其底部镶嵌了同底的正四棱锥形实心装饰块，容器内盛有 \(a\text{ 升}\) 水时，水面恰好经过正四棱锥的顶点 \(P\). 如果将容器倒置，水面也恰好过点 \(P\)（图 2）. 有下列四个命题：

\begin{circlelist}
\item 正四棱锥的高等于正四棱柱高的一半；
\item 将容器侧面水平放置时，水面也恰好过点 \(P\)；
\item 任意摆放该容器，当水面静止时，水面都恰好经过点 \(P\)；
\item 若往容器内再注入 \(a\text{ 升}\) 水，则容器恰好能装满.
\end{circlelist}
其中真命题的代号是\fillinblank{}.（写出所有真命题的代号）

\bitmapfigure[max width=0.55\linewidth]{img/2008/jiangxi_science/q16_fig1.png}
\end{problem}

\begin{answer}
B，D
\end{answer}

\begin{solution}
设正四棱柱的底面边长为 \(s\)，高为 \(H\)，正四棱锥的高为 \(h\). 在图 1 中，水面经过顶点 \(P\)，所以水的体积为正四棱柱内高为 \(h\) 的部分减去正四棱锥体积，即
\[
a=s^2h-\frac13s^2h=\frac23s^2h
\]
倒置后，点 \(P\) 距容器底面为 \(H-h\)，水面经过 \(P\)，且水面以下没有装饰块，所以
\[
a=s^2(H-h)
\]
两式相等，得
\[
\frac23h=H-h,\quad\text{即}\quad H=\frac53h
\]
因此正四棱锥的高不是正四棱柱高的一半，命题 A 错误.

容器除去实心装饰块后的总容积为
\[
V=s^2H-\frac13s^2h
=s^2\left(\frac53h-\frac13h\right)
=\frac43s^2h=2a
\]
所以再注入 \(a\) 升水后恰好装满，命题 D 正确.

当容器的一个侧面水平放置时，经过 \(P\) 且平行于该侧面的水平面是容器和正四棱锥的对称面，分别把容器空间和装饰块体积平分. 由 \(V=2a\)，当前水量正好是可用容积的一半，故水面恰好经过 \(P\)，命题 B 正确.

说明命题 C 错误：将正四棱锥的一个侧面水平放置，若水面经过 \(P\)，水面就必须与该侧面重合. 在原坐标系中取底面为
\(-\frac{s}{2}\leqslant x\)，\(y\leqslant\frac{s}{2}\)、\(z=0\)，顶面为 \(z=H\)，并令 \(P\left(0,0,h\right)\). 取包含底边 \(y=\frac{s}{2}\) 的侧面，其所在平面为

\[
y+\frac{s}{2h}z=\frac{s}{2}
\]

该平面一侧包含整个正四棱锥. 由 \(H=\frac53h\)，棱柱位于该侧的体积为

\[
V_-=\int_0^H s\left(s-\frac{s}{2h}z\right)\,\mathrm dz
=s^2\left(H-\frac{H^2}{4h}\right)=\frac{35}{36}s^2h
\]

因此该侧的可用空间体积为

\[
V_--\frac13s^2h=\frac{23}{36}s^2h
\]

另一侧的可用空间体积为

\[
V-\left(V_--\frac13s^2h\right)=\frac{25}{36}s^2h
\]

而当前水量 \(a=\frac23s^2h=\frac{24}{36}s^2h\)，两者均不等于 \(a\). 所以并非任意摆放时水面都经过 \(P\). 综上，真命题为 B，D.
\end{solution}

\section{解答题}

\begin{problem}
在 \(\triangle ABC\) 中，\(a\)，\(b\)，\(c\) 分别为角 \(A\)，\(B\)，\(C\) 所对的边长，\(a = 2\sqrt{3}\)，\(\tan \frac{A + B}{2} + \tan \frac{C}{2} = 4\)，\(\sin B \sin C = \cos^2 \frac{A}{2}\)，求 \(A\)，\(B\) 及 \(b\)，\(c\).
\end{problem}

\begin{solution}
因为 \(A+B+C=\pi\)，所以 \(\tan\frac{A+B}{2}=\cot\frac{C}{2}\). 令 \(t=\tan\frac{C}{2}\)，则 \(t>0\)，且
\(t+\dfrac{1}{t}=4\)，即 \(t^2-4t+1=0\)，从而 \(t=2-\sqrt{3}\) 或 \(t=2+\sqrt{3}\). 因此 \(C=30^\circ\) 或 \(C=150^\circ\).

又 \(A=\pi-B-C\)，由已知条件得 \(\sin B\sin C=\sin^2\dfrac{B+C}{2}\). 利用积化和差公式，左边为 \(\dfrac{\cos(B-C)-\cos(B+C)}{2}\)，右边为 \(\dfrac{1-\cos(B+C)}{2}\)，所以 \(\cos(B-C)=1\).

由 \(\cos(B-C)=1\) 且 \(B\)，\(C\in(0,\pi)\)，可知 \(B=C\). 若 \(C=150^\circ\)，则 \(B=150^\circ\)，从而 \(A+B+C>180^\circ\)，不合题意. 因此 \(B=C=30^\circ\)，\(A=120^\circ\). 由正弦定理，\(\dfrac{a}{\sin A}=\dfrac{b}{\sin B}=\dfrac{c}{\sin C}\)，且 \(\dfrac{a}{\sin A}=\dfrac{2\sqrt{3}}{\sqrt{3}/2}=4\). 所以 \(b=4\sin30^\circ=2\)，\(c=4\sin30^\circ=2\).
\end{solution}

\begin{problem}
因冰雪灾害，某柑桔基地果林严重受损，为此有关专家提出两种拯救果树的方案，每种方案都需分两年实施. 若实施方案一，预计第一年可以使柑桔产量恢复到灾前的 \(1.0\text{ 倍}\)，\(0.9\text{ 倍}\)，\(0.8\text{ 倍}\) 的概率分别是 \(0.3\)，\(0.3\)，\(0.4\)；第二年可以使柑桔产量为第一年产量的 \(1.25\text{ 倍}\)，\(1.0\text{ 倍}\) 的概率分别是 \(0.5\)，\(0.5\). 若实施方案二，预计第一年可以使柑桔产量达到灾前的 \(1.2\text{ 倍}\)，\(1.0\text{ 倍}\)，\(0.88\text{ 倍}\) 的概率分别是 \(0.2\)，\(0.3\)，\(0.5\)；第二年可以使柑桔产量为第一年产量的 \(1.2\text{ 倍}\)，\(1.0\text{ 倍}\) 的概率分别是 \(0.4\)，\(0.6\). 实施每种方案第一年与第二年相互独立，令 \( \xi_i \)（\(i = 1\)，\(2\)）表示方案 \(i\) 实施两年后柑桔产量达到灾前产量的倍数.
\begin{enumerate}
\item 写出 \(\xi_1\)，\(\xi_2\) 的分布列；
\item 实施哪种方案，两年后柑桔产量超过灾前产量的概率更大？
\item 不管哪种方案，如果实施两年后柑桔产量达不到，恰好达到，超过灾前产量，预计利润分别为 \(10\text{ 万元}\)，\(15\text{ 万元}\)，\(20\text{ 万元}\). 问实施哪种方案的平均利润更大？
\end{enumerate}
\end{problem}

\begin{solution}
\begin{enumerate}
\item 方案一中，两年后的产量倍数等于第一年倍数与第二年倍数的乘积. 因为两年相互独立，分别计算各组合的概率并合并相同取值，得到
\(P(\xi_1=0.8)=0.4\times0.5=0.2\)，\(P(\xi_1=0.9)=0.3\times0.5=0.15\)，\(P(\xi_1=1)=0.3\times0.5+0.8\times0.5=0.35\)，\(P(\xi_1=1.125)=0.3\times0.5=0.15\)，\(P(\xi_1=1.25)=0.3\times0.5=0.15\). 因此 \(\xi_1\) 的分布列为

\begin{tabular}{c|ccccc}
\(\xi_1\) & \(0.8\) & \(0.9\) & \(1\) & \(1.125\) & \(1.25\) \\
\hline
\(P\) & \(0.2\) & \(0.15\) & \(0.35\) & \(0.15\) & \(0.15\)
\end{tabular}

方案二中，可能的两年后产量倍数及概率为
\(P(\xi_2=0.88)=0.5\times0.6=0.3\)，\(P(\xi_2=1)=0.3\times0.6=0.18\)，\(P(\xi_2=1.056)=0.5\times0.4=0.2\)，\(P(\xi_2=1.2)=0.2\times0.6+0.3\times0.4=0.24\)，\(P(\xi_2=1.44)=0.2\times0.4=0.08\). 因此 \(\xi_2\) 的分布列为

\begin{tabular}{c|ccccc}
\(\xi_2\) & \(0.88\) & \(1\) & \(1.056\) & \(1.2\) & \(1.44\) \\
\hline
\(P\) & \(0.3\) & \(0.18\) & \(0.2\) & \(0.24\) & \(0.08\)
\end{tabular}

\item 方案一两年后产量超过灾前产量的概率为 \(P(\xi_1>1)=0.15+0.15=0.30\). 方案二两年后产量超过灾前产量的概率为 \(P(\xi_2>1)=0.20+0.24+0.08=0.52\). 因为 \(0.52>0.30\)，所以实施方案二的概率更大.

\item 方案一中，两年后产量达不到、恰好达到、超过灾前产量的概率分别为 \(0.2+0.15=0.35\)，\(0.35\)，\(0.30\)，所以平均利润为 \(E_1=10\times0.35+15\times0.35+20\times0.30=14.75\text{ 万元}\).

方案二中，上述三种情况的概率分别为 \(0.30\)，\(0.18\)，\(0.20+0.24+0.08=0.52\)，所以平均利润为 \(E_2=10\times0.30+15\times0.18+20\times0.52=16.1\text{ 万元}\). 因为 \(16.1>14.75\)，所以实施方案二的平均利润更大.
\end{enumerate}
\end{solution}

\begin{problem}
\begin{enumerate}
\item 等差数列 \(\{a_{n}\}\) 各项均为正整数，\(a_{1} = 3\). 前 \(n\) 项和为 \(S_{n}\)，等比数列 \(\{b_{n}\}\) 中，\(b_{1} = 1\)，且 \(b_{2}S_{2} = 64\)，\(\{b_{a_{n}}\}\) 是公比为 \(64\) 的等比数列. 求 \( a_{n} \) 与 \( b_{n} \)；
\item 证明：\(\frac{1}{S_1} + \frac{1}{S_2} + \cdots + \frac{1}{S_n} < \frac{3}{4}\).
\end{enumerate}
\end{problem}

\begin{solution}
\begin{enumerate}
\item 设等差数列 \(\{a_n\}\) 的公差为 \(d\)，等比数列 \(\{b_n\}\) 的公比为 \(q\). 因为 \(a_n=3+(n-1)d\) 对任意正整数 \(n\) 都是正整数，所以 \(d\) 是整数且 \(d\geqslant0\). 又 \(b_1=1\)，所以 \(b_n=q^{n-1}\)，并且 \(S_2=3+(3+d)=6+d\). 由 \(b_2S_2=64\) 得 \(q(6+d)=64\)，从而 \(q>0\).

由 \(b_{a_n}=q^{a_n-1}=q^{2+(n-1)d}\)，可知 \(\{b_{a_n}\}\) 的公比为 \(q^d\)，故 \(q^d=64\). 因为 \(d=0\) 时 \(q^d=1\)，不满足 \(q^d=64\)，所以 \(d>0\). 由 \(q=\dfrac{64}{d+6}\) 可知 \(q\) 是正有理数. 设 \(q=\dfrac{r}{s}\)，其中 \(r\)，\(s\) 是互质的正整数，则 \(r^d=64s^d\)，故 \(s^d\mid r^d\). 由 \(r\)，\(s\) 互质可得 \(s=1\)，所以 \(q\) 是正整数. 再由 \(q^d=64=2^6\)，正整数 \(q\) 只能为 \(2^k\)，且 \(kd=6\)，于是 \((q,d)\) 只能是 \((64,1)\)，\((8,2)\)，\((4,3)\) 或 \((2,6)\). 代入 \(q(6+d)=64\) 检验，只有 \(q=8\)，\(d=2\) 满足条件. 所以 \(a_n=3+2(n-1)=2n+1\)，\(b_n=8^{n-1}\).

\item 由 \(a_n=2n+1\) 得
\(S_n=\dfrac{n\left(a_1+a_n\right)}{2}=\dfrac{n(3+2n+1)}{2}=n(n+2)\). 因而
\(\dfrac{1}{S_n}=\dfrac{1}{n(n+2)}=\dfrac{1}{2}\left(\dfrac{1}{n}-\dfrac{1}{n+2}\right)\)，从而
\(\dfrac{1}{S_1}+\dfrac{1}{S_2}+\cdots+\dfrac{1}{S_n}=\dfrac{1}{2}\left(1+\dfrac{1}{2}-\dfrac{1}{n+1}-\dfrac{1}{n+2}\right)=\dfrac{3}{4}-\dfrac{2n+3}{2(n+1)(n+2)}<\dfrac{3}{4}\).
\end{enumerate}
\end{solution}

\begin{problem}
\begin{enumerate}
\item 如图，正三棱锥 \(O - ABC\) 的三条侧棱 \(OA\)，\(OB\)，\(OC\) 两两垂直，且长度均为 \(2\). \(E\)，\(F\) 分别是 \(AB\)，\(AC\) 的中点，\(H\) 是 \(EF\) 的中点，过 \(EF\) 的一个平面与侧棱 \(OA\)，\(OB\)，\(OC\) 或其延长线分别相交于 \(A_{1}\)，\(B_{1}\)，\(C_{1}\)，已知 \(OA_{1} = \frac{3}{2}\). 证明： \(B_{1}C_{1}\perp\) 平面 \(OAH\)；
\item 求二面角 \( O - A_{1}B_{1} - C_{1} \) 的大小.
\end{enumerate}
\bitmapfigure[max width=0.45\linewidth]{img/2008/jiangxi_science/q20_fig1.png}
\end{problem}

\begin{solution}
\begin{enumerate}
\item 以 \(OA\)，\(OB\)，\(OC\) 所在直线分别为 \(x\)，\(y\)，\(z\) 轴建立空间直角坐标系，则 \(O(0,0,0)\)，\(A(2,0,0)\)，\(B(0,2,0)\)，\(C(0,0,2)\). 由中点坐标公式，得 \(E(1,1,0)\)，\(F(1,0,1)\)，\(H\left(1,\dfrac{1}{2},\dfrac{1}{2}\right)\). 已知 \(OA_1=\dfrac{3}{2}\)，故 \(A_1\left(\dfrac{3}{2},0,0\right)\).

平面 \(A_1EF\) 的方程为 \(2x+y+z=3\)，所以它与 \(OB\)，\(OC\) 的交点分别为 \(B_1(0,3,0)\)，\(C_1(0,0,3)\). 因此 \(\overrightarrow{B_1C_1}=(0,-3,3)\)，\(\overrightarrow{OA}=(2,0,0)\)，\(\overrightarrow{AH}=\left(-1,\dfrac{1}{2},\dfrac{1}{2}\right)\). 于是 \(\overrightarrow{B_1C_1}\cdot\overrightarrow{OA}=0\)，\(\overrightarrow{B_1C_1}\cdot\overrightarrow{AH}=0\). 由于 \(\overrightarrow{OA}\) 与 \(\overrightarrow{AH}\) 不共线，\(\overrightarrow{B_1C_1}\) 是平面 \(OAH\) 的法向量方向. 又平面 \(OAH\) 的方程为 \(z-y=0\)，点 \(Q\left(0,\frac32,\frac32\right)\) 同时在直线 \(B_1C_1\) 和平面 \(OAH\) 上，因此 \(B_1C_1\perp\) 平面 \(OAH\).

\item 平面 \(OA_1B_1\) 的方程为 \(z=0\)，其一个法向量为 \(\bs{n}_1=(0,0,1)\). 平面 \(A_1B_1C_1\) 的方程为 \(2x+y+z=3\)，其一个法向量为 \(\bs{n}_2=(2,1,1)\). 设二面角 \(O-A_1B_1-C_1\) 的平面角为 \(\theta\)，取其锐角，则
\(\cos\theta=\dfrac{|\bs{n}_1\cdot\bs{n}_2|}{|\bs{n}_1||\bs{n}_2|}=\dfrac{1}{\sqrt{6}}\). 因此 \(\tan\theta=\sqrt{5}\)，且 \(0<\theta<\dfrac{\pi}{2}\)，所以二面角的大小为 \(\arctan\sqrt{5}\).
\end{enumerate}
\end{solution}

\begin{problem}
\begin{enumerate}
\item 设点 \(P(x_0, y_0)\) 在直线 \(x = m\)（\(y_0 \neq \pm m\)，\(0 < m < 1\)）上，过点 \(P\) 作双曲线 \(x^2 - y^2 = 1\) 的两条切线 \(PA\)，\(PB\). 切点为 \(A\)，\(B\)，定点 \(M\left(\frac{1}{m}, 0\right)\). 过点 \(A\) 作直线 \(x - y = 0\) 的垂线. 垂足为 \(N\)，试求 \(\triangle AMN\) 的重心 \(G\) 所在的曲线方程；
\item 求证：\(A\)，\(M\)，\(B\) 三点共线.
\end{enumerate}
\end{problem}

\begin{solution}
\begin{enumerate}
\item 设切点 \(A\left(x_1,y_1\right)\)，则 \(x_1^2-y_1^2=1\). 直线 \(x-y=0\) 的方向向量为 \((1,1)\)，故从 \(A\) 向该直线作垂线所得垂足为 \(N\left(\dfrac{x_1+y_1}{2},\dfrac{x_1+y_1}{2}\right)\). 设 \(\triangle AMN\) 的重心为 \(G(X,Y)\)，则
\(X=\dfrac{x_1+\frac{1}{m}+\frac{x_1+y_1}{2}}{3}=\dfrac{3x_1+y_1+\frac{2}{m}}{6}\)，\(Y=\dfrac{y_1+\frac{x_1+y_1}{2}}{3}=\dfrac{x_1+3y_1}{6}\). 从而 \(X-\dfrac{1}{3m}=\dfrac{3x_1+y_1}{6}\)，所以
\(\left(X-\dfrac{1}{3m}\right)^2-Y^2=\dfrac{(3x_1+y_1)^2-(x_1+3y_1)^2}{36}=\dfrac{8(x_1^2-y_1^2)}{36}=\dfrac{2}{9}\). 因此 \(G\) 所在的曲线方程为 \(\left(X-\dfrac{1}{3m}\right)^2-Y^2=\dfrac{2}{9}\)，将坐标变量仍记为 \(x\)，\(y\)，即为 \(\left(x-\dfrac{1}{3m}\right)^2-y^2=\dfrac{2}{9}\).

\item 设切点 \(A\left(x_1,y_1\right)\)，\(B\left(x_2,y_2\right)\). 双曲线 \(x^2-y^2=1\) 在点 \(\left(x_i,y_i\right)\) 处的切线方程为 \(x_i x-y_i y=1\). 因为点 \(P(m,y_0)\) 在两条切线上，所以
\(mx_1-y_0y_1=1\)，\(mx_2-y_0y_2=1\). 这说明 \(A\)，\(B\) 均在直线 \(mx-y_0y=1\) 上，而点 \(M\left(\dfrac{1}{m},0\right)\) 也满足 \(m\cdot\dfrac{1}{m}-y_0\cdot0=1\). 因此 \(A\)，\(M\)，\(B\) 三点共线.
\end{enumerate}
\end{solution}

\begin{problem}
\begin{enumerate}
\item 已知函数 \(f(x) = \frac{1}{\sqrt{1 + x}} + \frac{1}{\sqrt{1 + a}} + \sqrt{\frac{ax}{ax + 8}}\)，\(x \in (0, +\infty)\). 当 \(a = 8\) 时，求 \(f(x)\) 的单调区间；
\item 对任意正数 \(a\)，证明：\(1 < f(x) < 2\).
\end{enumerate}
\end{problem}

\begin{solution}
\begin{enumerate}
\item 当 \(a=8\) 时，
\(f(x)=\dfrac{1}{3}+\dfrac{1}{\sqrt{1+x}}+\sqrt{\dfrac{x}{1+x}}=\dfrac{1}{3}+\dfrac{1+\sqrt{x}}{\sqrt{1+x}}\). 求导得
\(f'(x)=\dfrac{1-\sqrt{x}}{2\sqrt{x}(1+x)^{3/2}}\). 因为分母在 \(x>0\) 时恒为正，所以当 \(0<x<1\) 时 \(f'(x)>0\)，当 \(x>1\) 时 \(f'(x)<0\). 因此 \(f(x)\) 在 \((0,1]\) 上单调递增，在 \([1,+\infty)\) 上单调递减.

\item 令 \(u=x\)，\(v=a\)，\(w=\dfrac{8}{ax}\)，则 \(u\)，\(v\)，\(w>0\)，且 \(uvw=8\)，并且
\(f(x)=\dfrac{1}{\sqrt{1+u}}+\dfrac{1}{\sqrt{1+v}}+\dfrac{1}{\sqrt{1+w}}\). 先证下界. 因为 \(u>0\)，有 \(\dfrac{1}{\sqrt{1+u}}>\dfrac{1}{1+u}\)，所以
\(f(x)>\sum\limits_{\mathrm{cyc}}\dfrac{1}{1+u}\). 另一方面，
\(\sum\limits_{\mathrm{cyc}}\dfrac{1}{1+u}-1=\dfrac{u+v+w-6}{(1+u)(1+v)(1+w)}\geqslant0\)，其中最后一步用了 \(u+v+w\geqslant3\sqrt[3]{uvw}=6\). 因此 \(f(x)>1\).

再证上界. 不妨设 \(u\leqslant v\leqslant w\). 若 \(u+v\geqslant6\)，则 \(v\geqslant3\)，\(w\geqslant3\)，从而 \(\dfrac{1}{\sqrt{1+v}}\leqslant\dfrac{1}{2}\)，\(\dfrac{1}{\sqrt{1+w}}\leqslant\dfrac{1}{2}\)，且 \(\dfrac{1}{\sqrt{1+u}}<1\)，故 \(f(x)<2\).

若 \(u+v<6\)，因为 \(u>0\)，有 \(2\sqrt{1+u}<u+2\)，从而
\(\dfrac{1}{\sqrt{1+u}}<\dfrac{u+2}{2(1+u)}=1-\dfrac{u}{2(1+u)}\). 对 \(v>0\) 同样有
\(2\sqrt{1+v}<v+2\)，故
\(\dfrac{1}{\sqrt{1+v}}<\dfrac{v+2}{2(1+v)}=1-\dfrac{v}{2(1+v)}\). 又由均值不等式，
\(\dfrac{u}{2(1+u)}+\dfrac{v}{2(1+v)}\geqslant\sqrt{\dfrac{uv}{(1+u)(1+v)}}\).

由于 \(u+v<6\)，有 \((1+u)(1+v)=uv+u+v+1<uv+8\). 再由 \(w=\dfrac{8}{uv}\)，得
\(\dfrac{uv}{(1+u)(1+v)}>\dfrac{uv}{uv+8}=\dfrac{1}{1+w}\)，从而 \(\dfrac{u}{2(1+u)}+\dfrac{v}{2(1+v)}>\dfrac{1}{\sqrt{1+w}}\). 因此
\(f(x)<2-\dfrac{u}{2(1+u)}-\dfrac{v}{2(1+v)}+\dfrac{1}{\sqrt{1+w}}<2\). 综上，对任意正数 \(a\)，恒有 \(1<f(x)<2\).
\end{enumerate}
\end{solution}
